\section{Conclusion}\label{sec:conclusion}

This study re-examined the efficiency of Arctic shipping by isolating the effect of routing methodology. Using Rotterdam--Yokohama as a representative corridor, five case studies compared great-circle routing, sea-only graph routing, fuel and emission estimates, auxiliary loads, and operational sensitivity. The results demonstrate that assumptions about routing methodology can fundamentally alter conclusions: while \ac{gc} routing suggests the \ac{nsr} is up to 40--50\% shorter than the \ac{suez} route, a sea-only routing approach constrained by land masking produces the opposite outcome, with the Arctic route emerging longer and more emission-intensive.

These findings challenge the prevailing narrative that the Arctic offers inherent efficiency gains. By quantifying how methodological bias propagates through distance, time, fuel, and emissions, the study shows that Arctic advantages are not universal but conditional on modeling choices. This has direct implications for sustainability assessments and policy debates, particularly as the shipping sector seeks alignment with \ac{imo} decarbonization targets.

The analysis also highlighted the importance of including auxiliary loads and operational choices such as speed, which exert strong influence on total fuel use and emissions. These factors amplify differences between routes and emphasize that evaluations cannot be based on geometry alone. 

Future work (Paper~3 onward) will build on this methodological baseline by integrating dynamic environmental constraints, including sea ice concentration, wind, and wave conditions. By progressively adding realism, the research will enable a more accurate appraisal of when, if ever, Arctic routes can be competitive with established traditional passages.
